% ===================================================================
%  Minimalbeispiel: Liedblatt – „Grün, grün, grün sind alle meine Kleider“
%  (Volkslied, gemeinfrei), G-Dur, 2/4-Takt. Erste Strophe unter den Noten
%  (musixlyr), Strophen 2–6 zweispaltig darunter, alles auf einer Seite.
%  Eigenständig: braucht keine header.tex.
%
%  Drei Durchläufe:  pdflatex … → musixflx … → pdflatex …
%  (musixflx gleicht alle Notenzeilen auf dieselbe Länge aus.)
% ===================================================================
\documentclass[a4paper,12pt]{article}
\usepackage[T1]{fontenc}
\usepackage[utf8]{inputenc}
\usepackage[ngerman]{babel}
\usepackage{lmodern}
\usepackage[margin=1.8cm]{geometry}
\usepackage{multicol}       % Strophen zweispaltig
\usepackage{musixtex}
\input musixlyr             % Liedtext unter den Noten
\pagestyle{empty}
\setlength{\columnsep}{1em}

% Eine Strophe: \strophe{Nummer}{Text}; bleibt beim Spaltenumbruch zusammen.
\newcommand{\strophe}[2]{%
  \noindent\makebox[1.4em][l]{\textbf{#1.}}%
  \begin{minipage}[t]{\dimexpr\linewidth-1.4em}\raggedright #2\end{minipage}%
  \par\vspace{1.2em}}

\begin{document}
\begin{center}
{\LARGE\bfseries\sffamily Grün, grün, grün sind alle meine Kleider}\\[4pt]
{\small Text und Melodie: Volksweise \textbullet\ G-Dur}
\end{center}
\vspace{1em}

\Largemusicsize             % großes System: VOR \begin{music}
\begin{music}
\parindent0pt                % erste Zeile nicht einrücken: alle Zeilen gleich lang
\staffbotmarg=8\Interligne  % Platz unter jedem System (innerhalb von music)
\instrumentnumber{1}\setclef1\treble
\generalmeter{\meterfrac24}\generalsignature{1}\nobarnumbers
\setsongraise1{10mm}        % Text dicht unter die eigene Zeile
\setlyrics{strophe}{Grün, ja grün sind al-le mei-ne Klei-der,
grün, ja grün ist al-les, was ich hab.
Da-rum lieb ich al-les, was so grün ist,
weil mein Schatz ein Jä-ger ist.}
\assignlyrics1{strophe}
\startpiece
% Tonhöhen: d = d', e = e', f = fis' (Tonart), g = g', h = a', i = h', j = c'', k = d''.
% Zwei Noten auf einer Silbe (Melisma): \lyricsoff vor der zweiten Note,
% \lyricson danach; dazu ein Bindebogen \islurd … \tslur (oben: \isluru).
\NOtes\qu g\qu g\en\bar                                              % Grün, ja
\NOtes\qu h\qu h\en\bar                                              % grün sind
\Notes\ibu0f0\qb0f\tbu0\qb0f\ibu0e1\qb0e\tbu0\qb0f\en\bar            % al-le mei-ne
\NOtes\qu g\qu d\en                                                  % Klei-der,
\alaligne
\NOtes\qu g\qu g\en\bar                                              % grün, ja
\NOtes\qu h\qu h\en\bar                                              % grün ist
\Notes\ibu0f0\qb0f\tbu0\qb0f\ibu0e1\qb0e\tbu0\qb0f\en\bar            % al-les, was ich
\NOtes\qu g\qp\en                                                    % hab.
\alaligne
\NOtes\ql i\en\Notes\ibl0i2\isluru0i\qb0i\lyricsoff\tslur0k\tbl0\qb0k\lyricson\en\bar % Da-rum
\NOtes\ql j\ql j\en\bar                                              % lieb ich
\Notes\ibu0h0\qb0h\tbu0\qb0h\ibu0h2\qb0h\tbu0\qb0j\en\bar            % al-les, was so
\NOtes\ql i\ql i\en                                                  % grün ist,
\alaligne
\NOtes\qu g\qu g\en\bar                                              % weil mein
\NOtes\qu h\qu h\en\bar                                              % Schatz ein
\Notes\ibu0f0\islurd0f\qb0f\lyricsoff\tslur0f\tbu0\qb0f\lyricson     % Jä-
\ibu0e1\islurd0e\qb0e\lyricsoff\tslur0f\tbu0\qb0f\lyricson\en\bar    % ger
\NOtes\qu g\qp\en                                                    % ist.
\Endpiece
\end{music}
\assignlyrics1{}            % Zuweisung beenden (gilt sonst global weiter)

\vspace{2em}

\begin{multicols}{2}
\strophe{2}{Blau, ja blau sind alle meine Kleider,\\
blau, ja blau ist alles, was ich hab.\\
Darum lieb ich alles, was so blau ist,\\
weil mein Schatz ein Seemann ist.}

\strophe{3}{Weiß, ja weiß sind alle meine Kleider,\\
weiß, ja weiß ist alles, was ich hab.\\
Darum lieb ich alles, was so weiß ist,\\
weil mein Schatz ein Bäcker ist.}

\strophe{4}{Schwarz, ja schwarz sind alle meine Kleider,\\
schwarz, ja schwarz ist alles, was ich hab.\\
Darum lieb ich alles, was so schwarz ist,\\
weil mein Schatz ein Schornsteinfeger ist.}

\strophe{5}{Rot, ja rot sind alle meine Kleider,\\
rot, ja rot ist alles, was ich hab.\\
Darum lieb ich alles, was so rot ist,\\
weil mein Schatz ein Reiter ist.}

\strophe{6}{Bunt, ja bunt sind alle meine Kleider,\\
bunt, ja bunt ist alles, was ich hab.\\
Darum lieb ich alles, was so bunt ist,\\
weil mein Schatz ein Maler ist.}
\end{multicols}

\end{document}
